% =====================================================================
% docs/metricas/metricas.tex -- Bitacora de metricas de implementacion
% Proyecto II CE-1113: reproductor de audio en DE1-SoC (P2, integracion)
%
% Una fila por compilacion relevante (fin de cada fase). Los datos salen
% de hw/quartus/output_files/: de1_soc.flow.rpt (Flow Summary),
% de1_soc.fit.rpt (Fitter RAM Summary, recursos por entidad) y
% de1_soc.sta.rpt (Setup/Hold/Fmax). Sirve de base para T8.1-T8.3.
% Compilar: pdflatex metricas.tex
% =====================================================================
\documentclass[11pt]{article}
\usepackage[utf8]{inputenc}
\usepackage[T1]{fontenc}
\usepackage[spanish,es-noshorthands]{babel}
\usepackage[margin=2.2cm]{geometry}
\usepackage{booktabs}
\usepackage{siunitx}
\sisetup{group-separator={\,}, group-minimum-digits=4, output-decimal-marker={.}}
\usepackage{array}
\usepackage[hidelinks]{hyperref}

\title{Métricas de implementación del SoC\\\large Proyecto II CE-1113 — Reproductor de audio en DE1-SoC}
\author{P2 — Integración del SoC}
\date{Última actualización: 28 de noviembre de 2026}

\begin{document}
\maketitle

\section*{Condiciones}
Quartus Prime Lite 25.1std.0 Build 1129, dispositivo Cyclone V 5CSEMA5F31C6
(\num{32070} ALMs, \num{4065280} bits de memoria en bloque, 397 M10K, 87 DSP).
Reloj del sistema: \texttt{CLOCK\_50}, restringido a \SI{20}{\nano\second} en
\texttt{hw/quartus/de1\_soc.sdc}.

\section{Recursos por fase (sistema completo)}

\begin{table}[h]
\centering
\small
\begin{tabular}{@{}l l r r r r r r@{}}
\toprule
Fase & Commit & ALMs & Registros & Bits de memoria & M10K & DSP & Pines \\
\midrule
1 — Nios V/m + RAM 128 KiB + JTAG UART & \texttt{0eae00a}
  & \num{2279} (7\,\%) & \num{3113} & \num{1051648} (26\,\%) & 132 (33\,\%) & 0 & 15 \\
\bottomrule
\end{tabular}
\caption{Flow Summary al cierre de cada fase.}
\end{table}

\section{Tiempos (Timing Analyzer)}

\begin{table}[h]
\centering
\small
\begin{tabular}{@{}l l r r r r@{}}
\toprule
Fase & Reloj & Setup slack & Hold slack & Fmax & Modelo peor caso \\
\midrule
1 & \texttt{CLOCK\_50}            & \SI{8.528}{\nano\second}  & \SI{0.279}{\nano\second} & \SI{87.17}{\mega\hertz} & Slow 1100mV 85C \\
1 & \texttt{altera\_reserved\_tck} & \SI{11.460}{\nano\second} & \SI{0.374}{\nano\second} & \SI{96.03}{\mega\hertz} & Slow 1100mV 85C \\
\bottomrule
\end{tabular}
\caption{Slack positivo $\Rightarrow$ el diseño cumple. Fmax es la frecuencia
máxima de la ruta más lenta: $1/(20 - 8.528)\,\text{ns} \approx \SI{87.2}{\mega\hertz}$.}
\end{table}

\section{Desglose de la memoria en bloque (fase 1)}

\begin{table}[h]
\centering
\small
\begin{tabular}{@{}l l r@{}}
\toprule
Memoria & Organización & Bits \\
\midrule
\texttt{nios\_ram} (código, datos, pila) & $32768 \times 32$ & \num{1048576} \\
Banco de registros del Nios V (2 copias, 2 puertos de lectura) & $2 \times (32 \times 32)$ & \num{2048} \\
FIFO de TX y RX de la JTAG UART & $2 \times (64 \times 8)$ & \num{1024} \\
\midrule
Total (coincide con el Flow Summary) & & \num{1051648} \\
\bottomrule
\end{tabular}
\caption{Del \emph{Fitter RAM Summary}. La ROM y los CSR del módulo de depuración
del Nios van en MLAB (ALMs usados como memoria) y no cuentan aquí.}
\end{table}

\section{Referencia: módulo propio \texttt{audio\_dsp} aislado}
Medido por P1 (commit \texttt{f5784ff}, \texttt{docs/verifica/recursos\_audio\_dsp.txt}),
con los puertos como pines virtuales: \num{937} ALMs, \num{1796} registros,
512 bits de memoria (1 M10K) y 3 DSP. Se repetirá dentro del sistema completo
en la fase 2 para T8.1.

\section{Observaciones}
\begin{itemize}
  \item Fase 1: la RAM del Nios es el 99.7\,\% de la memoria usada. El presupuesto
        restante (\num{3013632} bits, unos 368 KiB) debe cubrir la RAM compartida
        de 64 KiB, las FIFO de audio y el buffer del eco.
\end{itemize}

\end{document}
